% Study sheet: Volume I, Chapter 3, ML Lifecycle.
% Numbers: \val{...} from ml_lifecycle_studysheet_numbers.tex, written by ml_lifecycle_studysheet_numbers.py.
\documentclass{studysheet}
\input{ml_lifecycle_studysheet_numbers}
\sheetvolume{1}
\sheetchapter{3}{ML Lifecycle}
\sheetsources{Numbers from the chapter's screening-clinic scenarios in mlsysim; assumed inputs are marked.}

\begin{document}
% ===================================================================== FRONT: LEARN
\frontside
\begin{bigidea}
An ML system is built in six stages that \hl{hand each other artifacts}: targets, a dataset, weights,
release evidence, a running service. Because its behavior is learned from data, the world can change
it after launch, so the stages form a \hl{loop}. A limit found late costs every artifact built without it.
\end{bigidea}

\begin{leftcol}
\sheetsec{Key ideas}
\begin{ideas}
\idea{The lifecycle is a loop.}{Unlike rule-based software, a model can degrade with unchanged code when
  its inputs change, so monitoring feeds evidence back.}
\idea{Constraints propagate both ways.}{Device limits reach back to model and data. Found late, a limit
  forces repair of every artifact built without it.}
\idea{A request is not a requirement.}{State a measurable gain over a named baseline, split into
  statistical, physical and operational limits, set before training.}
\idea{Contracts at every hand-off.}{Each stage states its inputs, outputs and a quality check. A fix
  starts at the earliest stage whose contract missed the condition.}
\idea{Evaluation measures; validation decides.}{Curated accuracy is not readiness. A training run
  counts only if it can be reproduced.}
\idea{Iteration speed buys search breadth.}{Cycle time bounds how many ideas a team can test; it does
  not guarantee accuracy.}
\end{ideas}

\vspace{11pt}
\sheetsec{Six stages, and what each hands on}
{\centering
\begin{tikzpicture}[sheet, stg/.style={accbox, text width=46pt, align=center, minimum height=22pt}]
  \node[stg] (s1) at (30,66) {1 Problem definition};
  \node[stg] (s2) at (125,66) {2 Data collection};
  \node[stg] (s3) at (220,66) {3 Model training};
  \node[stg] (s4) at (220,0) {4 Evaluation, validation};
  \node[stg] (s5) at (125,0) {5 Deployment};
  \node[stg, hotbox] (s6) at (30,0) {6 Monitoring};
  \draw[arr] (s1) -- node[note, above=2pt]{targets} (s2);
  \draw[arr] (s2) -- node[note, above=2pt]{dataset} (s3);
  \draw[arr] (s3) -- node[note, left=3pt]{weights} (s4);
  \draw[arr] (s4) -- node[note, below=2pt]{manifest} (s5);
  \draw[arr] (s5) -- node[note, below=2pt]{service} (s6);
  \draw[arr, accent, dashed] (s6) -- node[note, right=3pt]{evidence} (s1);
  \node[small, anchor=south] at (s2.north) [yshift=2pt] {Data};
  \node[small, anchor=south] at (s3.north) [yshift=2pt] {Algorithm};
  \node[small, anchor=north] at (s5.south) [yshift=-2pt] {Machine};
  \node[note, align=center] at (125,33) {a fix starts at the\\earliest stage whose\\contract missed it};
\end{tikzpicture}\par}

\vspace{11pt}
\sheetsec{Where the hours go}
\begin{tiles}
\tile{\val{datawork}\%}{of data scientists: cleaning or collecting data takes the most time}\hfill
\tile{hours vs.\ months}{model cycle: hours to a week; data cycle: weeks to months}
\end{tiles}
\end{leftcol}%
\begin{rightcol}
\sheetsec{The release gate}
\begin{sheettable}{@{}l L@{}}
\thead{Check} & \thead{Passes when}\\ \midrule
Floors & one threshold gives sensitivity (sick patients flagged) $>$ \val{sensfloor}\% and specificity (healthy cleared) $>$ \val{specfloor}\%\\
Subgroups & every slice (age, camera, cataract) clears the floors\\
Calibration & stated confidence matches how often it is right\\
Device & fits the \val{budget} MB reservation; latency met on the clinic device\\
Field & holds across sites, blurred or dim images, and months\\
\end{sheettable}
\vspace{3pt}
{\fontsize{8.6}{11}\selectfont\color{muted}Accuracy (correct $\div$ all) is dominated by the healthy
majority, hence two floors. A failed row holds the
release.\par}

\vspace{6pt}
\fcard{Upload time}%
  {\Time = \frac{\DataVolume}{\Bandwidth}}%
  {Bytes $\div$ link rate (8 bits per byte). Busy clinic: \val{cpat} $\times$ \val{cphoto} photos $\times$
   \val{cmb} MB $=$ \val{cdaily} MB at \val{cmbit} Mb/s $=$ \val{crate} MB/s: \val{csec} s (\val{chours} h),
   \val{cshare}\% of an \val{cshift}-hour shift (\val{cshiftsec} s); one photo \val{cone} s.}%
  {Results are due within seconds, link or not, so the model runs in the clinic and sends
   a \val{csumkb} KB summary per patient, \val{cred}$\times$ less.}{}

\vspace{11pt}
\sheetsec{Monitoring: three clocks}
\begin{sheettable}{@{}L l L@{}}
\thead{Signal} & \thead{Arrives} & \thead{Tells you}\\ \midrule
Latency, errors, queue & minutes & the service runs\\
Input drift, confidence, referral and rejection rates & hours to days & inputs or behavior changed\\
Graded outcomes: sensitivity per slice & weeks & accuracy changed\\
\end{sheettable}
\vspace{3pt}
{\fontsize{8.6}{11}\selectfont\color{muted}Drift warns; labels measure. Only referred cases get
graded, so grade a sample of cleared ones too. An alert opens a review under approved exposure limits.\par}
\end{rightcol}

% ===================================================================== BACK: PRACTICE
\backside
\begin{answers}
\ans{Finish}{Collect \val{ycollect} + grade \val{ygrade} + prepare \val{wprep} + train \val{wtrain} +
  validate \val{wval} + package \val{wpack} $=$ \val{yrepair} days. Existing images, labels and preparation
  code survive; the manifest records the cataract slice against the floor. (Contracts; release gate)}
\ans{1}{\val{qpat} $\times$ \val{cphoto} $\times$ \val{cmb} $=$ \val{qdaily} MB $\div$ \val{qrate} MB/s $=$
  \val{qsec} s (\val{qhours} h), \val{qshare}\% of \val{cshiftsec} s. One photo takes \val{qone} s, too slow for
  a seconds deadline: run the model in the clinic. (Upload time)}
\ans{2}{Accuracy \val{alwaysacc}\%, sensitivity zero, specificity \val{alwaysspec}\%. Requirement: find more
  referable disease than today's referral path, sensitivity $>$ \val{sensfloor}\% and specificity
  $>$ \val{specfloor}\% in every subgroup, within device limits. (Requirement idea; gate)}
\columnbreak
\ans{3}{Unknown: drift shows changed inputs, not accuracy, and cleared patients go ungraded. Grade a
  sample of cleared and referred images; meanwhile, run under approved exposure limits. If
  confirmed, fix at the owning stage and re-pass the gate.}
\ans{4}{(a) Data collection: targeted images. (b) Deployment: the serving path, not the model.
  (c) Model training: record environment and seed per run. (Contracts; Evaluation)}
\end{answers}

\begin{leftcol}
\sheetsec{Worked example}
\begin{practicebox}\boxtitle{Day \val{dayfound}: the model needs \val{needgb} GB; the device reserves \val{budget} MB}
Handed off on day \val{dayhand} at \val{sensday}\% sensitivity, the model is \val{over}$\times$ too big. Who owns the fix? What does it cost?\par\vspace{5pt}
{\centering
\begin{tikzpicture}[sheet]
  \pgfmathsetmacro\u{200/\val{wrepair}}
  \pgfmathsetmacro\xa{40+\val{wprep}*\u}  \pgfmathsetmacro\xb{\xa+\val{wtrain}*\u}
  \pgfmathsetmacro\xc{\xb+\val{wval}*\u}  \pgfmathsetmacro\xd{\xc+\val{wpack}*\u}
  \node[note, anchor=east] at (36,8) {repair};
  \draw[accbox] (40,0) rectangle (\xa,16);  \node[lab] at ({(40+\xa)/2},8) {prepare};
  \draw[accbox] (\xa,0) rectangle (\xb,16); \node[lab] at ({(\xa+\xb)/2},8) {train};
  \draw[accbox] (\xb,0) rectangle (\xc,16); \node[lab] at ({(\xb+\xc)/2},8) {validate};
  \draw[accbox] (\xc,0) rectangle (\xd,16); \node[lab] at ({(\xc+\xd)/2},8) {package};
  \foreach \d in {0,...,6}{\draw[dim] ({40+\d*\u},0) -- ({40+\d*\u},-2.5);
    \node[small, anchor=north] at ({40+\d*\u},-4) {\d};}
  \node[small, anchor=north east] at (\xd,-12) {days};
  \node[note, anchor=south west] at (40,19) {kept: raw images and labels};
\end{tikzpicture}\par}\vspace{4pt}
\begin{steps}
\item Walk back: evaluation tested only sensitivity, training got no memory budget, problem definition
  never recorded \val{budget} MB. \res{The fix starts at problem definition.}
\item Redo preparation, weights, validation and package, in sequence: \val{wprep} + \val{wtrain} + \val{wval}
  + \val{wpack} $=$ \res{\val{wrepair} days}.
\item Lost beyond repair: the part of the \val{builtdays} days spent tuning a model too large to fit.
\end{steps}
\takeaway{Put the device limit in the problem definition.}
\end{practicebox}

\vspace{12pt}
\sheetsec{Your turn: finish it}
\begin{practicebox}\boxtitle{A subgroup fails the release gate}
Sensitivity is \val{sensall}\% overall, \val{senscat}\% for cataract patients (floor \val{sensfloor}\%).
Provenance shows few cataract images.\par\vspace{3pt}
\textbf{Given:} release held; the owner is data collection.
\assumed{} collecting takes \val{ycollect} days, grading \val{ygrade}; other tasks as above.\par
\textbf{You:} list what to redo, in order, the elapsed days, what survives, and what the manifest records.
\end{practicebox}
\end{leftcol}%
\begin{rightcol}
\sheetsec{Common mistakes}
\mistake{``It matched specialists on curated images, so it is ready for clinics.''}%
  {In Thai clinics, \val{thairej} of \val{thaisub} photos (\val{thaipct}\%) failed the quality check.
   Validation must sample field conditions.}{}
\mistake{``The bigger model starts \val{startgap} points ahead, so it wins.''}%
  {\val{slowgain} vs.\ \val{fastgain} points per run, ceiling \val{ceiling}\%. In \val{weeks} weeks, weekly
   runs from \val{slowstart}\% give \val{slowruns} experiments (\val{slowfinal}\%); hourly runs from
   \val{faststart}\% allow \val{fastruns} useful ones (\val{fastfinal}\%): about the same accuracy, \val{ideas}$\times$ the ideas tested.}{}
\mistake{``We fixed the random seed, so training is reproducible.''}%
  {Parallel sums arrive in varying order and floating-point addition is not associative, so
   weights differ. Record data version, code and environment; measure the run-to-run spread.}{}

\sheetsec{Check yourself}
\begin{qs}
\item \assumed{} a clinic sees \val{qpat} patients a day, \val{cphoto} photos of \val{cmb} MB each, on a
  \val{qmbit} Mb/s uplink and an \val{cshift}-hour shift. Raw upload time, and share of the shift? Can a cloud result meet the seconds deadline?
\item \assumed{} \val{prev}\% of screened patients have referable disease. A model always answers
  ``healthy.'' Its accuracy, sensitivity and specificity? Rewrite
  ``detect retinopathy with high accuracy'' as a requirement.
\item A clinic swaps its camera. Next day the input-drift alert fires for that site; referral rate and
  latency are unchanged. Has sensitivity dropped? What evidence would tell, and what do you do meanwhile?
\item Which stage owns each fix? (a) One clinic's new camera gives unfamiliar images. (b) Serving resizes images differently
  from training. (c) Two engineers on the same data version get different sensitivity; neither recorded environment or seed.
\end{qs}
\end{rightcol}
\end{document}
